\section{Certified finite-chain separations}\label{app:xx84-certificate}

This appendix proves Theorems~\ref{thm:xx84-separation}
and~\ref{thm:xx-assisted-separation}. We first give the chain-only
certificate and then the assisted certificate at smaller sizes.
\paragraph{Exact cross covariance and dyadic quadrature.}
Put $M=N+1$ and $m=N/4$. The exact covariance in the proof of
Corollary~\ref{cor:xx-mesh} has two cross-cut parity blocks. On the
even-even block, multiplying rows and columns by signs gives the symmetric
matrix with entries
\begin{equation}\label{eq:xx84-kernel}
 B_{ij}=\frac{1}{2M}\left[
 \csc\!\left(\frac{\pi(2i+2j-1)}{2M}\right)+
 \sec\!\left(\frac{\pi(2i-2j)}{2M}\right)\right],
 \qquad 1\le i,j\le m.
\end{equation}
The two parity blocks have identical singular-value lists. In fact the
left covariance $C_A$ satisfies $\Gamma C_A\Gamma=I-C_A$, where $\Gamma$
alternates signs between the sites. This follows from the sign reversal of
the hopping Hamiltonian under $\Gamma$ and the absence of a zero energy
at even length. Thus in parity order $C_A-I/2$ has off-diagonal block $D$
and its adjoint, with $D$ square. The product of the full covariance
cross block and its adjoint equals $C_A-C_A^2$, whose parity blocks are
$I/4-DD^*$ and $I/4-D^*D$.

Partition the index interval into the integer bins
$[a,\min(2a-1,m)]$ for $a=1,2,4,\ldots,m$. On a bin of $L$ points use
Gaussian quadrature of order $k=\min(18,L)$ for counting measure.
For a bin with midpoint $b$ the nodes are $b+L\mu_i$, where $\mu_i$
are eigenvalues of the $k\times k$ Jacobi matrix with diagonal zero
and off-diagonal entries
\[
 \beta_j=\frac{j}{2L}\sqrt{\frac{L^2-j^2}{4j^2-1}},\qquad 1\le j<k.
\]
The weights are $L|v_i(1)|^2$ for unit eigenvectors $v_i$. These
recurrence coefficients follow by taking the ratios of the squared
norms of monic orthogonal polynomials on $0,\ldots,L-1$. For degree $j$
the monic polynomial is the $j$th forward difference of
$x^{\underline j}(x-L)^{\underline j}$ divided by $(2j)!/j!$.
Summation by parts gives orthogonality to lower degrees and squared
norm $(j!)^4\binom{L+j}{2j+1}/(2j)!$ by Vandermonde's identity.
The symmetric grid gives the zero diagonal after centering. The
quadrature moments agree through degree $2k-1$ because a path from
the first basis vector back to itself in that many Jacobi steps cannot
use any index beyond $k$. The nodes lie in the bin since the Jacobi
matrix compresses multiplication by the centered grid coordinate.
At $k=L$ the nodes and weights are the
original integers and unit weights.

Form the exact quadrature matrix $K$ by evaluating~\eqref{eq:xx84-kernel}
at the nodes and multiplying each entry by the geometric mean of the
associated weights. Write $s(A)$ for the decreasing singular values of
a matrix $A$, padding by zeros when comparing lists. We show
\begin{equation}\label{eq:xx84-compression}
 \|s(B)-s(K)\|_{\ell^2}<6\times10^{-6}.
\end{equation}
For a pair of bins, decompose the kernel into
$1/(\pi w)+M^{-1}[f(w/M)+g(z/M)]$ with $w=2i+2j-1$, $z=2i-2j$,
$f(t)=(2\sin(\pi t/2))^{-1}-(\pi t)^{-1}$, and
$g(t)=(2\cos(\pi t/2))^{-1}$. Approximate each of the three terms by
its degree-$17$ Taylor polynomial about the two bin midpoints. These
polynomials in the bin variables have separate degrees at most 17.
The corresponding weighted quadrature matrix therefore has the same
nonzero singular values as the polynomial matrix evaluated on the full
integer grid. This follows from the isometric evaluation of polynomials
on each bin by exact quadrature for their pairwise inner products.

For the first term, $|w-w_0|/w_0\le1/3$ on a dyadic bin pair, so the
geometric series error is at most $3^{-18}/(\pi w)$. On the original
grid, the Frobenius norm of $1/(\pi w)$ is at most
$\sqrt{1+\ln(2m)}/(2\pi)<4/3$ by summing over $i+j-1$. The corresponding
quadrature Frobenius norm is at most twice as large, since the maximum
of $w$ on a bin pair is less than twice its minimum and the quadrature
weights sum to the number of grid points in each bin.
On $|t|=3/2$, $|f(t)|<1$. Indeed, either $|\Im t|\ge1/2$ or
$\sqrt2\le |\Re t|\le 3/2$, and each case gives
$|\sin(\pi t/2)|\ge1/\sqrt2$. On $|t|=3/4$ we have $|g(t)|\le4/3$
since $|\cos(\pi t/2)|\ge\cos(3\pi/8)>3/8$. These bounds hold
in the enclosed disks by analyticity. Write $t_0$ for a bin-pair center
and $\delta$ for the argument's maximum deviation on the bin rectangle.
The bin lengths give $\delta\le1/4$. The argument interval is symmetric
about $t_0$, so $\delta\le R_0-|t_0|$, with $R_0=1$ for $f$ and
$R_0=1/2$ for $g$. Consequently
$\delta/(R-|t_0|)\le(1/4)/(R-R_0+1/4)$, where $R=3/2$ and $3/4$,
respectively. These ratios are at most $1/3$ and $1/2$.
The Cauchy estimates bound the two Taylor
remainders by $(3/2)3^{-18}$ and $(8/3)2^{-18}$. Their contribution to
the sum of the two Frobenius errors is at most their sum times
$2m/M<1/2$. The full sum is less than $6\times10^{-6}$. The
singular-value perturbation inequality in Frobenius norm gives
\eqref{eq:xx84-compression}.

\paragraph{Verification of the compressed covariance.}
The code \texttt{certify\_xx84\_modes.py} uses binary64 eigensolvers to
propose eigenpairs. It checks them by the following enclosures. Write
$\widehat K$ for the symmetric binary64 matrix produced by the code.
For round-to-nearest arithmetic put $u_0=2^{-53}$ and
$\gamma_j=ju_0/(1-ju_0)$. A product of $j$ relative rounding factors
or their inverses differs from one by at most $\gamma_j$.

\begin{enumerate}
\item \emph{Nodes and weights.} Rational square comparisons verify each
Jacobi entry to within $2\times10^{-15}$, so the operator-norm perturbation
is at most $4\times10^{-15}$. For each proposed pair $(\mu,v)$, exact
rational checks give $|\|v\|^2-1|<2\times10^{-14}$ and rounded-matrix
residual norm below $2\times10^{-14}$. The normalized residual on the
exact matrix is bounded by
\[
 \frac{2\times10^{-14}}{\sqrt{1-2\times10^{-14}}}+4\times10^{-15}
 < r_J,\qquad r_J=4\times10^{-14}.
\]
The computed gaps exceed $10^{-3}$, so the disjoint intervals of radius
$r_J$ each contain a distinct exact eigenvalue and exhaust the spectrum.
Expanding $v/\|v\|$ in the exact eigenvectors bounds the norm of its
component perpendicular to the matching eigenvector by
$r_J/(10^{-3}-r_J)$. After sign alignment and restoring the original norm,
the vector error is bounded by
$\sqrt2\,r_J/(10^{-3}-r_J)+2\times10^{-14}<b_v=5.66\times10^{-11}$.
Since the checked first-coordinate square of $v$ exceeds $10^{-3}$, the
relative weight error after the two multiplications is at most
\[
 \frac{2b_v+b_v^2+\gamma_2(1+2\times10^{-14})}{10^{-3}-2b_v-b_v^2}
 < \epsilon_w,
 \qquad \epsilon_w=2\times10^{-7}.
\]
The bin start $a$ and length $L$ are powers of two, and $L\le a$.
In forming $a+(L-1)/2+L\mu$, the multiplication by $L$ is exact and
$|\mu|\le1/2+r_J$. Thus the node's absolute error is bounded by
\[
 Lr_J+a\gamma_2+\tfrac L2\gamma_3+L(\tfrac12+r_J)u_0
 < a\epsilon_x,\qquad \epsilon_x=5\times10^{-14}.
\]
The exact node is at least $a$, giving a relative error below $\epsilon_x$.
The integer nodes and unit weights in the short bins are represented exactly.

\item \emph{Kernel entries.} The Machin enclosure verifies an absolute
error below $10^{-15}$ for the binary64 value of $\pi$, hence relative
error below $\epsilon_\pi=10^{-15}/3$. Rounding $H=2M$ to $2N$ gives
relative error $1/M<u_0$. For nodes $x,y$ let $t_+=\pi(2x+2y-1)/H$
and $t_-=\pi|2x-2y|/H$. Substituting the rounded nodes into $2x+2y-1$
contributes relative error at most $4\epsilon_x/3$ since $2x+2y-1\ge3$.
The rounded nodes are also at least one and their doubling is exact.
The addition and subtraction in the rounded numerator give relative error
at most $(4/3)u_0+u_0(1+(4/3)u_0)<3u_0$. For the difference argument,
substituting the nodes contributes absolute error at most
$(\pi/2)\epsilon_x$, since $x+y\le2m$. With the remaining rounding factors
and the enclosure of $\pi$, these estimates give
\begin{align*}
 \left|\frac{\widehat t_+}{t_+}-1\right|
 &\le (1+4\epsilon_x/3)(1+3u_0)(1+\epsilon_\pi)(1+\gamma_3)-1
 <3\times10^{-13},\\
 |\widehat t_--t_-|
 &\le \tfrac\pi2\epsilon_x+
     (\tfrac\pi4+\tfrac\pi2\epsilon_x)
     \bigl[(1+\epsilon_\pi)(1+\gamma_4)-1\bigr]
 <3\times10^{-13}.
\end{align*}
The rounded arguments satisfy $|\widehat t_+|\le1.58$ and
$|\widehat t_-|\le0.8$. Their nonzero magnitudes exceed $2^{-137}$,
since doubled rounded nodes are multiples of $2^{-51}$ and the rounded
denominator is $2^{85}$. Thus the argument squares are normal when nonzero.
For the degree-20 even Taylor polynomials the first omitted terms are
below $10^{-17}$ and the sums of absolute terms are below $3$ on these
intervals. The ratios of adjacent term magnitudes are below $0.42$,
so each Horner sum retains at least half the magnitude of its coefficient.
The smallest coefficient has magnitude above $2^{-66}$. This also keeps
nonzero Horner products in the normal range. Each term passes through
at most 31 rounding factors, counting coefficient rounding, the powers
of the rounded square, and the Horner multiplications and additions.
The polynomial evaluation and truncation error is therefore bounded by
$3\gamma_{31}+10^{-17}<10^{-13}$ for each function.

At the exact arguments, $\operatorname{sinc}(t_+)\ge2/\pi$ and
$\cos(t_-)\ge1/\sqrt2$. The function
$\operatorname{sinc}(t)=\int_0^1\cos(zt)\,dz$ has Lipschitz constant
$1/2$, and cosine has Lipschitz constant one. With $b=3\times10^{-13}$,
the relative errors of the evaluated sine and cosine are bounded by
\begin{align*}
 &(1+b)\bigl[1+(\pi/2)(10^{-13}+(\pi/4)b)\bigr](1+u_0)-1
   <10^{-12},\\
 &\sqrt2(10^{-13}+b)<10^{-12},
\end{align*}
respectively. The rounded reciprocal of either factor then has relative
error at most $(10^{-12}+u_0)/(1-10^{-12})<2\times10^{-12}$.

Let $p$ be the product of the two exact weights, $\widehat p$ their
rounded counterpart, and $h$ the proposed square root of $\widehat p$.
The rounded weights lie between $2^{-10}$ and $2^{83}$. The rounded
square and division in the root check produce a quotient within
$10^{-14}$ of one. The final subtraction of one is exact by Sterbenz's
lemma. The check forces the rounded square and quotient into the normal
range, so allowing for those two roundings gives
\[
 \left|\frac{h^2}{\widehat p}-1\right|
 \le(1+10^{-14})(1+\gamma_2)-1<1.1\times10^{-14}.
\]
The weight enclosures then put $h^2/p$ between
$(1-\epsilon_w)^2(1-u_0)(1-1.1\times10^{-14})$ and
$(1+\epsilon_w)^2(1+u_0)(1+1.1\times10^{-14})$. These endpoints lie
strictly between $(1-3\times10^{-7})^2$ and $(1+3\times10^{-7})^2$.
The two reciprocal terms are positive, so their sum has the same
relative error bound before rounding. Including the addition, denominator
rounding, division and final multiplication gives
\[
 \left|\frac{\widehat K_{ij}}{K_{ij}}-1\right|
 \le (1+2\times10^{-12})(1+3\times10^{-7})(1+\gamma_4)-1
 <10^{-6}.
\]

\item \emph{Frobenius norm and eigenpairs.} The entry calculation uses
normal numbers or zero throughout. In particular its positive entries
exceed $2^{-100}$ by the weight bounds and the upper bounds on the
evaluated sine and cosine. The rounded sum of squared entries is below
four. With $d<1600$ the matrix dimension, squaring and the two sums
use at most $2d+1$ rounding factors per term. The absolute error is below
$4\gamma_{2d+1}/(1-\gamma_{2d+1})<3\times10^{-12}$. Hence
\[
 \|K-\widehat K\|_F
 \le \frac{10^{-6}}{1-10^{-6}}\|\widehat K\|_F
 < \frac{10^{-6}\sqrt{4+3\times10^{-12}}}{1-10^{-6}}
 <2.1\times10^{-6}.
\]
For each proposed leading eigenpair $(\mu,v)$ the code verifies
$|\mu|<1$, a computed squared norm within $10^{-11}$ of one, and
each computed residual component below $10^{-14}$ in absolute value.
The positive-sum error bound gives $\|v\|>0.99$. The arithmetic error in
the matrix-vector product is at most $\gamma_{d+1}\|\widehat K\|_F\|v\|$
apart from underflow. Allowing $2^{-1075}$ per scalar operation for gradual
underflow, its total contribution after summation, taking norms and
normalizing is below $2^{-1000}$. The scalar multiplication contributes
at most $u_0\|v\|$ apart from underflow, and the final subtraction costs
a factor $1/(1-u_0)$ in recovering the norm before rounding. Thus
\[
 \frac{\|\widehat K v-\mu v\|}{\|v\|}
 < \frac{40\times10^{-14}}{0.99(1-u_0)}
    +2.01\gamma_{d+1}+u_0+2^{-1000}<3\times10^{-12}.
\]
The selected signed intervals with this radius are disjoint and each
contains an exact eigenvalue. Subtracting the sum of the squared lower
bounds on their absolute values from the upper bound on the squared
Frobenius norm controls the remaining squared singular mass.
\end{enumerate}

For the finite model take amplitudes with single-mode ratios
$r_i=(32767/32768)^{k_i}$, where the 81 integers $k_i$ are listed in Table~\ref{tab:xx84-model}. Each mode occurs twice. The model
singular values are $s_i^*=r_i/(1+r_i^2)$. The code encloses each
$r_i$ by binary exponentiation with directed decimal rounding and
compares $s_i^*$ to the interval for the assigned exact singular value.
The reference run bounds the omitted squared mass and the squared error
of the retained singular values by $3.46\times10^{-11}$ and
$4.73\times10^{-11}$. Match each selected exact singular value to its
assigned $s_i^*$ and every other singular value to zero. Sorting both
lists in decreasing order cannot increase their $\ell^2$ distance,
by the rearrangement inequality. The verifier checks that the two squared
errors sum to less than $(4\times10^{-5})^2$. Combining this bound with
the two matrix approximation errors gives one-parity singular-list
distance at most $4.81\times10^{-5}$. The verifier also checks that the
largest model singular value plus this distance is below $0.499$.


\begin{table}[htb]
\centering\footnotesize
\begin{tabular}{@{}rrrrrrrrr@{}}
\toprule
2676 & 8030 & 13389 & 18758 & 24138 & 29530 & 34936 & 40355 & 45787 \\
51231 & 56686 & 62153 & 67629 & 73114 & 78607 & 84109 & 89617 & 95132 \\
100653 & 106180 & 111713 & 117250 & 122793 & 128340 & 133892 & 139448 & 145007 \\
150571 & 156139 & 161710 & 167284 & 172862 & 178443 & 184027 & 189614 & 195204 \\
200797 & 206393 & 211991 & 217592 & 223195 & 228801 & 234410 & 240021 & 245634 \\
251249 & 256867 & 262486 & 268108 & 273732 & 279358 & 284986 & 290616 & 296248 \\
301882 & 307517 & 313155 & 318794 & 324435 & 330078 & 335722 & 341368 & 347016 \\
352665 & 358316 & 363968 & 369622 & 375278 & 380935 & 386594 & 392254 & 397915 \\
403578 & 409242 & 414908 & 420575 & 426243 & 431913 & 437584 & 443256 & 448929 \\
\bottomrule
\end{tabular}
\caption{The 81 exponents of the finite geometric model, read row by row.
Each exponent occurs twice in the Schmidt-mode product. The common
amplitude ratio is $32767/32768$.}\label{tab:xx84-model}
\end{table}

\paragraph{State and game loss.}
The Schmidt angle of a mode with singular value $s$ is
$\theta(s)=\tfrac12\arcsin(2s)$. The occupation and its complement are
the squares of the two Schmidt amplitudes. The derivative of $\theta$
is less than 16 on $0\le s\le0.499$. Choose the finite model state
$\widehat\psi$ in the exact Schmidt mode bases of the chain and freeze
the omitted modes. The two parity lists agree and each mode overlap is
the cosine of the angle difference. Hence
\[
 \|\psi_N-\widehat\psi\|^2
 =2\left(1-\prod_i\cos^2(\theta(s_i)-\theta(s_i^*))\right)
 \le 2\sum_i(\theta(s_i)-\theta(s_i^*))^2,
\]
which gives vector distance less than
$16\sqrt2\,(4.81\times10^{-5})<0.0011063<3/2500$.

\paragraph{Scoring the geometric model.}
Let $q=32767/32768$, and let $k_i$ be the integer exponents in
Table~\ref{tab:xx84-model}, one for each of the 81 retained mode pairs.
The model uses two modes with Schmidt ratio $q^{k_i}$ for every $i$. Put
$p_0=\prod_i(1+q^{2k_i})^{-2}$ and let $n_e$ be the coefficient of $z^e$ in
$\prod_i(1+z^{k_i})^2$. The Schmidt amplitudes are $\sqrt{p_0}q^e$, with
multiplicity $n_e$ in increasing order of $e$. Retain the prefix $e\le380000$.
For its unnormalised amplitudes $\lambda$, put
$\tau=1-\|\lambda\|_2^2$, $E=\|(I-D)\lambda\|_2^2$, and
$v=\|(I-P)\lambda\|_2^2$, where $P$ averages consecutive groups of nine
amplitudes and the prefix is padded with zeros. The players project locally
onto the corresponding Schmidt-prefix subspaces, which gives a common
branch on $\widehat\psi$. Apply~\eqref{eq:operational-adaptive} to the
prefix by homogeneity and use the classical strategy on the complement,
with loss contribution $\tau/27$. This gives
\[
 L_{\rm model}\le
 \frac{2(\sqrt E+\sqrt{8v/9})^2+72v+9\tau}{243}.
\]
The integer recurrence for $n_e$ adds $2n_{e-k_i}+n_{e-2k_i}$ for each
mode pair, processing $e$ in decreasing order and using zero for a negative
index. The prefix has rank $185{,}926{,}736$. The score verifier uses this run
representation throughout. For $E$, it intersects intervals of equal
amplitudes with the intervals obtained by doubling their endpoints. The
overlap on each intersection is the product of the corresponding two
amplitudes, its integer length, and $1/\sqrt2$. To compute $v$ it counts
only the nine-element groups that cross an interval endpoint or the end of
the prefix. Their contribution is the sum of the nine squared amplitudes
minus the square of their sum divided by nine.

Every interval endpoint in the verifier is an integer multiple of
$10^{-45}$. Addition is exact on this grid. Products and quotients use
integer floor and ceiling, and square roots use integer square roots of the
endpoints multiplied by $10^{45}$. This gives
$\tau<0.000082218$, $E<0.013735797$, $v<0.000800859$, and
$L_{\rm model}<0.000410720762$. Use the same local POVMs on $\psi_N$.
Their average losing effect $F$ satisfies $0\le F\le I$, so
$\sqrt{L_{\rm XX}}=\|F^{1/2}\psi_N\|
\le\|F^{1/2}\widehat\psi\|+\|\psi_N-\widehat\psi\|$.
With the aligned vector-distance bound $3/2500$, the physical loss obeys
\[
 L_{\rm XX}\le(\sqrt{L_{\rm model}}+3/2500)^2.
\]
Using the directed upper endpoint for $L_{\rm model}$, the integer
score verifier \texttt{model\_score.py} gives the upper endpoint
\[
 0.000460799754443163771129846265247324951980458
 <0.000461<\frac{3-2\sqrt2}{243}.
\]

